% !Mode:: "TeX:UTF-8"
% !TEX program  = xelatex
% !BIB program  = biber

\documentclass[a4paper,UTF8,scheme=chinese,zihao=-4,twoside]{ctexart}

% 基本宏包
\usepackage{geometry}
\usepackage{fontspec}
\usepackage{titlesec}
\usepackage{setspace}
\usepackage{tocloft}
\usepackage{fancyhdr}
\usepackage{graphicx}
\usepackage{amsmath,amssymb,amsthm}
\usepackage{enumitem}
\usepackage{hyperref}
\usepackage{caption}
\usepackage{booktabs}
\usepackage{subcaption}
% \usepackage{float}  % 注释掉，和 floatrow 冲突
\usepackage{listings}
\usepackage{xcolor}
\usepackage{etoolbox}
\usepackage[style=gb7714-2015,gbpunctin=false]{biblatex}

% 参考文献资源文件
\addbibresource{references.bib}

% 页面设置
\geometry{a4paper,top=27mm,bottom=31mm,left=25mm,right=25mm,bindingoffset=5mm}

% 字体设置 - 严格按照原模板
% 英文使用 Times New Roman
\IfFileExists{fonts/times/times.ttf}{
    \setmainfont[
        Path=fonts/times/,
        Extension = .ttf,
        UprightFont = *,
        BoldFont = *bd,
        ItalicFont = *i,
        BoldItalicFont = *bi
    ]{times}
}{
    \IfFontExistsTF{Times New Roman}{
        \setmainfont{Times New Roman}
    }{
        \IfFontExistsTF{TeX Gyre Termes}{
            \setmainfont{TeX Gyre Termes}
        }{}
    }
}

% 中文使用宋体
\IfFileExists{fonts/simsun.ttc}{
    \setCJKmainfont[
        Path=fonts/,
        AutoFakeBold
    ]{simsun.ttc}
    \setCJKsansfont[
        Path=fonts/,
        AutoFakeBold
    ]{simhei.ttf}
}{
    \setCJKmainfont{SimSun}[AutoFakeBold]
    \setCJKsansfont{SimHei}[AutoFakeBold]
}

% 字号命令
\newcommand{\初号}{\zihao{0}}
\newcommand{\小初}{\zihao{-0}}
\newcommand{\一号}{\zihao{1}}
\newcommand{\小一}{\zihao{-1}}
\newcommand{\二号}{\zihao{2}}
\newcommand{\小二}{\zihao{-2}}
\newcommand{\三号}{\zihao{3}}
\newcommand{\小三}{\zihao{-3}}
\newcommand{\四号}{\zihao{4}}
\newcommand{\小四}{\zihao{-4}}
\newcommand{\五号}{\zihao{5}}
\newcommand{\小五}{\zihao{-5}}
\newcommand{\六号}{\zihao{6}}
\newcommand{\小六}{\zihao{-6}}
\newcommand{\七号}{\zihao{7}}
\newcommand{\八号}{\zihao{8}}

% 字体命令
\newcommand{\宋体}{\songti}
\newcommand{\黑体}{\heiti}
\newcommand{\仿宋}{\fangsong}
\newcommand{\楷书}{\kaishu}
\newcommand{\粗体}{\bfseries}
\newcommand{\斜体}{\itshape}

% 章节标题格式
\ctexset{
    section = {
        name={{},.},
        format={\黑体\三号},
    },
    subsection = {
        format={\黑体\四号},
    },
    subsubsection = {
        format={\黑体\小四},
    },
}

% 章节标题间距
\titlespacing{\section}{0pt}{15.6pt}{.3ex plus .0ex}
\titlespacing{\subsection}{0pt}{1.2ex plus .0ex minus .0ex}{.3ex plus .0ex}

% 行距设置
\renewcommand{\baselinestretch}{1.72}
\usepackage[MSWordLineSpacingMultiple=1.5]{zhlineskip}

% 页眉页脚设置
\pagestyle{fancy}
\lhead{}\chead{}\rhead{}\lfoot{}\rfoot{}
\cfoot{\五号\thepage}
\renewcommand{\headrulewidth}{0bp}
\renewcommand{\footrulewidth}{0bp}

% 图表标题格式
\DeclareCaptionLabelSeparator{enspace}{\enspace}
\DeclareCaptionLabelFormat{rightBracket}{#2)}
\renewcommand{\captionfont}{\bfseries\songti\五号}
\captionsetup[subfigure]{labelformat=rightBracket, labelsep=enspace}
\captionsetup[table]{labelsep=enspace}
\captionsetup[figure]{labelsep=enspace}
\ctexset{
    figurename = {图},
    tablename = {表},
    listfigurename = {插图},
    listtablename = {表格},
    proofname = {证明},
}

% 浮动体设置
\usepackage{floatrow}
\DeclareFloatFont{floatfont}{\五号}
\floatsetup[table]{
    capposition = top,
    objectset = centering,
    margins = centering,
    font = floatfont,
}
\floatsetup[figure]{
    capposition = bottom,
    objectset = centering,
    margins = centering,
    font = floatfont,
}

% 目录格式
\renewcommand\cfttoctitlefont{\hfill\黑体\小二}
\renewcommand\cftaftertoctitle{\hfill}
\renewcommand\cftsecfont{\粗体\宋体\四号}
\renewcommand\cftsecpagefont{\粗体\宋体\四号}
\renewcommand\cftsecleader{\cftdotfill{\cftdotsep}}
\renewcommand\cftsubsecfont{\宋体\四号}
\renewcommand\cftsubsecpagefont{\宋体\四号}
\renewcommand\cftsubsubsecfont{\宋体\四号}
\renewcommand\cftsubsubsecpagefont{\宋体\四号}
\setlength\cftparskip{7pt}
\setlength\cftbeforesecskip{0bp}
\setlength\cftsubsecindent{0bp}
\setlength\cftsubsubsecindent{0bp}
\setcounter{tocdepth}{3}

% 超链接设置
\hypersetup{
    breaklinks,
    hidelinks,
    hypertexnames=true,
    bookmarks=true,
    bookmarksopen=true,
    bookmarksnumbered=true,
    pdflang=zh-CN,
}

% 下划线命令（用于封面）- 简化版，不用 @
\newcommand\下划线[2][6em]{\underline{\makebox[#1][c]{#2}}}

% ========== 论文信息填写区 ==========
\newcommand{\tTitle}{南方科技大学毕业论文题目}
\newcommand{\tTitleEn}{Undergraduate Thesis Title}
\newcommand{\tSubtitle}{副标题（如无则留空）}
\newcommand{\tSubtitleEn}{Subtitle (optional)}
\newcommand{\tAuthor}{你的姓名}
\newcommand{\tStudentID}{12345678}
\newcommand{\tDepartment}{某某系}
\newcommand{\tDepartmentEn}{Department of XXX}
\newcommand{\tMajor}{某某专业}
\newcommand{\tMajorEn}{XXX Program}
\newcommand{\tAdvisor}{指导教师姓名}
\newcommand{\tAdvisorEn}{Advisor Name}
\newcommand{\tKeywords}{关键词1；关键词2；关键词3}
\newcommand{\tKeywordsEn}{Keyword1, Keyword2, Keyword3}
% =====================================

\begin{document}

% 中文封面
\thispagestyle{empty}
\newgeometry{top=2.25cm,bottom=2.25cm,left=2.54cm,right=2.48cm}
\begin{center}
    \begin{tabular}{llr}
        分类号 \下划线[2.3cm]{} \hspace{0.45\textwidth} & 编号\下划线[2.3cm]{} \\[15pt]
        U\hfill D\hfill C\下划线[2.3cm]{} \hspace{0.45\textwidth} & 密级\下划线[2.3cm]{}
    \end{tabular}
    \vskip 1.05cm
    \vskip 1.05cm
    {\粗体\宋体\小初 本科生毕业设计（论文）}
    \vskip 2.25cm
    {\粗体\宋体\三号
    \begin{tabular}{rl}
        题目： & \下划线[9cm]{\tTitle} \\[6pt]
        \ifdefempty{\tSubtitle}{}{\underline{\makebox[9cm][c]{\tSubtitle}} \\[20pt]}
        \ifdefempty{\tSubtitle}{\\[20pt]}{}
姓名：& \下划线[9cm]{\tAuthor} \\[16pt]
        学号： & \下划线[9cm]{\tStudentID} \\[16pt]
        院系： & \下划线[9cm]{\tDepartment} \\[16pt]
        专业： & \下划线[9cm]{\tMajor} \\[16pt]
        指导教师： & \下划线[9cm]{\tAdvisor} \\[16pt]
    \end{tabular}
}
\end{center}
\restoregeometry
\clearpage

% 英文封面
\thispagestyle{empty}
\newgeometry{top=2.54cm,bottom=2.75cm,left=2.54cm,right=2.48cm}
\begin{center}
    CLC \underline{\makebox[2.3cm][c]{}} \hfill Number\underline{\makebox[2.3cm][c]{}} \\[5pt]
    UDC\underline{\makebox[2.3cm][c]{}} \hfill Available for reference~\quad$\square$Yes~\quad$\square$No
    \vskip 1.8cm
    \vskip 1.6cm
    {\小初 Undergraduate Thesis}
    \vskip 2.7cm
    {\粗体\三号
    \begin{tabular}{rl}
        Thesis Title: & \underline{\makebox[9cm][c]{\tTitleEn}} \\[12pt]
        \ifdefempty{\tSubtitleEn}{}{\underline{\makebox[9cm][c]{\tSubtitleEn}} \\[15pt]}
        \ifdefempty{\tSubtitleEn}{\\[15pt]}{}
        Student Name: & \underline{\makebox[9cm][c]{\tAuthor}} \\[12pt]
        Student ID: & \underline{\makebox[9cm][c]{\tStudentID}} \\[12pt]
        Department: & \underline{\makebox[9cm][c]{\tDepartmentEn}} \\[12pt]
        Program: & \underline{\makebox[9cm][c]{\tMajorEn}} \\[12pt]
        Thesis Advisor: & \underline{\makebox[9cm][c]{\tAdvisorEn}} \\[12pt]
    \end{tabular}
}
\end{center}
\restoregeometry
\clearpage

% 诚信承诺书
\thispagestyle{empty}
\newgeometry{top=2.74cm,bottom=2.54cm,left=3.17cm,right=3.17cm}
\begin{center}
    {\bfseries\黑体\二号 诚信承诺}
\end{center}
\vskip 1.65cm
\begin{spacing}{1.85}\宋体\四号
    1.本人郑重承诺所呈交的毕业设计（论文），是在导师的指导下，独立进行研究工作所取得的成果，所有数据、图片资料均真实可靠。\par
    2.除文中已经注明引用的内容外，本论文不包含任何其他人或集体已经发表或撰写过的作品或成果。
    对本论文的研究作出重要贡献的个人和集体，均已在文中以明确的方式标明。\par
    3.本人承诺在毕业论文（设计）选题和研究内容过程中没有抄袭他人研究成果和伪造相关数据等行为。\par
    4.在毕业论文（设计）中对侵犯任何方面知识产权的行为，由本人承担相应的法律责任。

    \vskip 4.55cm
    \begin{flushright}
        \makebox[11em][l]{作者签名：} \\
        \makebox[11em][l]{\underline{\makebox[2cm][c]{}}年\underline{\makebox[1cm][c]{}}月\underline{\makebox[1cm][c]{}}日}
    \end{flushright}
\end{spacing}
\restoregeometry
\clearpage

% 罗马数字页码
\pagenumbering{Roman}
\setcounter{page}{1}

% 摘要标题页
\begin{center}
    \黑体\二号 \tTitle
    \\[12pt]
    \ifdefempty{\tSubtitle}{}{\hspace*{\fill} \黑体\小二 ——\tSubtitle \\}
    \vskip 1cm
    \宋体\四号 \tAuthor
    \\[9pt]
    （\楷书\小四{\tDepartment \quad 指导教师： \tAdvisor}）
\end{center}
\vskip 1.5cm

% 中文摘要
\begin{spacing}{1.45}
    \noindent\黑体\三号 [摘要]：
    \宋体\四号
    在此处输入中文摘要内容。摘要应概括论文的目的、方法、结果和结论，语言精炼准确。
    摘要通常不分段，顶格开始书写。

    \vfill
    \noindent\黑体\三号 [关键词]：
    \宋体\四号 \tKeywords
\end{spacing}
\clearpage

% 英文摘要
\begin{center}
    \二号 \tTitleEn
\end{center}
\vskip 1.5cm
\begin{spacing}{1.45}
    \noindent\三号 [ABSTRACT]：
    \四号
    Enter your English abstract content here. The abstract should summarize the purpose, methods, results, and conclusions of the thesis.

    \vfill
    \noindent\三号 [Keywords]：
    \四号 \tKeywordsEn
\end{spacing}
\clearpage

% 目录
\pdfbookmark[1]{目录}{toc}
\tableofcontents
\clearpage

% 正文开始，阿拉伯数字页码
\pagenumbering{arabic}
\setcounter{page}{1}

\section{引言}
在此处输入引言内容。介绍研究背景、目的和意义。

\section{相关工作}
在此处介绍相关领域的研究现状。可以使用~\cite{knuth1984}引用参考文献。

\section{方法}
在此处介绍你的研究方法。

\subsection{算法设计}
介绍算法的详细设计。

\subsection{实现细节}
介绍实现过程中的关键细节。

\section{实验结果与分析}
在此处展示实验结果并进行分析。

\subsection{实验设置}
介绍实验的设置和环境。

\subsection{结果分析}
分析实验结果。

\section{结论}
总结你的工作，并展望未来的研究方向。

% 参考文献
\clearpage
\phantomsection
\addcontentsline{toc}{section}{参考文献}
\section*{\centerline{参考文献}}
\printbibliography[heading=none]

% 附录（可选）
\clearpage
\phantomsection
\appendix
\addcontentsline{toc}{section}{附录}
\section*{\centerline{附录}}
\section{附录标题}
附录内容。

% 致谢
\clearpage
\phantomsection
\addcontentsline{toc}{section}{致谢}
\section*{\centerline{致谢}}
在此处感谢指导老师、同学和家人的支持与帮助。

\end{document}
